% chapters/03/checklist.tex
\bagian{Checklist Pencapaian Kompetensi}

Bab ini dinilai melalui unjuk kerja, bukan melalui hafalan. Centang setiap
butir berikut hanya apabila Anda benar-benar telah melakukannya pada komputer
Anda sendiri.

\begin{checklistbox}
\begin{checklist}
  \item Skrip pemeriksaan lingkungan berjalan dan melaporkan bahwa lingkungan
        siap dipakai.
  \item Saya dapat menyebutkan nama berkas penafsir Python yang dipakai
        komputer saya dan menjelaskan mengapa keterangan itu penting.
  \item Saya dapat memasang Pygame sendiri, termasuk ketika pemasangan gagal
        karena masalah izin.
  \item Saya dapat menjelaskan perbedaan kedudukan Pygame dan SDL.
  \item Saya dapat membuka jendela Pygame dengan ukuran yang saya tentukan
        sendiri melalui baris perintah.
  \item Saya dapat mengubah warna latar program dan melihat perubahannya pada
        jendela.
  \item Saya dapat menambahkan baris teks baru beserta posisinya.
  \item Saya dapat menjelaskan pekerjaan yang dilakukan
        \code{pygame.init()}.
  \item Saya dapat menjelaskan perbedaan menggambar dan menayangkan, serta
        pekerjaan \code{pygame.display.flip()}.
  \item Saya dapat menjelaskan perbedaan perulangan peristiwa pada bab ini
        dengan siklus permainan yang akan dibahas pada Bab~IV.
  \item Saya dapat membaca pesan kesalahan Python dan menentukan baris yang
        menjadi penyebabnya.
  \item Saya dapat membedakan peringatan dari kesalahan dan bersikap tepat
        terhadap keduanya.
  \item Repositori Git saya memuat folder \code{code}, \code{catatan}, dan
        \code{aset}, berkas \code{.gitignore}, serta penyimpanan awal dengan
        pesan yang menerangkan isinya.
  \item Saya telah menulis catatan praktikum yang memuat kendala dan cara
        penyelesaiannya.
\end{checklist}
\end{checklistbox}
